% area: Representation Theory

\begin{definition}[Group Representation]
\label{def:representation}
A \textit{representation} of a group $G$ on a vector space $V$ over a field $k$
is a group homomorphism
\[
  \rho : G \to GL(V).
\]
The dimension $\dim_k V$ is called the \textit{degree} of the representation.
\end{definition}

\begin{definition}[Subrepresentation]
\label{def:subrepresentation}
A subspace $W \subseteq V$ is a \textit{subrepresentation} (or $G$-invariant
subspace) if $\rho(g)(W) \subseteq W$ for all $g \in G$.
This notion relies on the setup of \ref{def:representation}.
\end{definition}

\begin{definition}[Irreducible Representation]
\label{def:irreducible}
A representation $(\rho, V)$ (in the sense of \ref{def:representation}) is
\textit{irreducible} (or \textit{simple}) if $V \neq 0$ and the only
subrepresentations (see \ref{def:subrepresentation}) of $V$ are $\{0\}$ and
$V$ itself.
\end{definition}

\begin{theorem}[Schur's Lemma]
\label{thm:schur}
Let $(\rho, V)$ and $(\sigma, W)$ be irreducible representations
(\ref{def:irreducible}) of $G$ over an algebraically closed field $k$.
\begin{enumerate}
  \item Any $G$-module homomorphism $\phi : V \to W$ is either zero or an
        isomorphism.
  \item $\mathrm{End}_G(V) \cong k$.
\end{enumerate}
\end{theorem}

\begin{proof}
(1) Since $\ker(\phi)$ and $\mathrm{Im}(\phi)$ are subrepresentations of $V$
and $W$ respectively, irreducibility forces each to be $\{0\}$ or the whole
space. The non-trivial cases give an isomorphism.

(2) For any $\lambda \in k$, the map $\phi - \lambda \cdot \mathrm{id}_V$ is
again a $G$-endomorphism. Since $k$ is algebraically closed, $\phi$ has some
eigenvalue $\lambda_0$, and $\phi - \lambda_0 \cdot \mathrm{id}_V$ has a
nontrivial kernel. By (1) it must be identically zero, so $\phi = \lambda_0
\cdot \mathrm{id}_V$.
\end{proof}

\begin{theorem}[Maschke's Theorem]
\label{thm:maschke}
Let $G$ be a finite group and $k$ a field with $\mathrm{char}(k) \nmid |G|$.
Then every representation (\ref{def:representation}) of $G$ over $k$ is
completely reducible, i.e.\ decomposes into irreducibles (\ref{def:irreducible}).
\end{theorem}

\begin{example}[Regular Representation]
\label{ex:regular}
The \textit{regular representation} of $G$ is the action of $G$ on the group
algebra $k[G]$ by left multiplication (a special case of \ref{def:representation}):
\[
  \rho(g) \cdot \sum_{h \in G} a_h \, h \;=\; \sum_{h \in G} a_h \, (gh).
\]
By \ref{thm:maschke}, over an algebraically closed field of characteristic zero,
every irreducible representation appears in $k[G]$ with multiplicity equal to
its degree.
\end{example}

\begin{remark}
\label{rem:conj-classes}
Over $\mathbb{C}$, the number of (isomorphism classes of) irreducible
representations of $G$ equals the number of conjugacy classes of $G$.
This follows from \ref{thm:maschke} and \ref{thm:schur}.
\end{remark}

\begin{corollary}[Dimension Formula]
\label{cor:dimension}
For a finite group $G$ over $\mathbb{C}$, combining \ref{thm:maschke} and
\ref{rem:conj-classes},
\[
  \sum_{\chi \in \hat{G}} (\dim V_\chi)^2 = |G|.
\]
\end{corollary}

\begin{problem}[Character Table of $S_3$]
\label{prob:s3}
Using \ref{cor:dimension} and \ref{thm:schur}, compute the complete character
table of the symmetric group $S_3$.
Verify the column orthogonality relations:
\[
  \sum_{\chi \in \hat{G}} \chi(g)\,\overline{\chi(h)}
  \;=\; |C_G(g)| \cdot \delta_{[g],[h]}.
\]
\end{problem}
